\section{Ernst Junger on Apoliteia}

The following passages are from Ernst Junger's philosophical-metahistorical novel \emph{Eumeswil}, a book which complements Julius Evola's \emph{Ride the Tiger} quite nicely.

\begin{quotex}
I am an anarch -- not because I despise authority, but because I need it. Likewise, I am not a nonbeliever, but a man who demands something worth believing in. ... Thus I take my duties seriously within an overall context that I reject for its mediocrity. The important thing is that my rejection actually refers to the totality and does not take up within it a stance that can be defined as conservative, reactionary, liberal, ironic, or in any way social.
\end{quotex}


\begin{quotex}
The lack of ideas or -- put more simply -- of gods causes an inexplicable moroseness, almost like a fog that the sun fails to penetrate. The world turns colorless; words lose substance, especially when they are to transcend sheer communication. ... Loss of history and decay of language are mutual determinants.
\end{quotex}


\begin{quotex}
Any man who swears allegiance to a political change is a fool, a \emph{facchino} for services that are not his business. The most rudimentary step toward freedom is to free oneself from all that. Basically each person senses it, and yet he keeps voting.''
\end{quotex}


\begin{quotex}
If and how far a mind penetrates matter, and whether it grasps the crown of the root from which the details branch off -- these things are perceived even in practice.
\end{quotex}


\begin{quotex}
Man should not be the sun's friend, but the sun itself. And that he \emph{is}; the mistake lies in his failure to recognize his place, his home, and thereby his right.
\end{quotex}



\flrightit{Posted on 2010-08-23 by Will}

\begin{center}* * *\end{center}

\begin{footnotesize}
\begin{sffamily}
\texttt{Simon Friedrich on 2010-09-08 at 12:38 said:}

Thanks Gornahoor -- always happy when someone else appreciates Jünger´s anarch!

Simon


\hfill

\end{sffamily}
\end{footnotesize}

